% Provenance (never printed): sources/book2/BRIEFS.md, g11-01. Marginal quantities as derivatives, the rule
% marginal benefit = marginal cost, average and marginal, tax rates recalled from Book 1 (g8), sunk costs
% recalled from Book 1 (g7). Real numbers: figdata/grade-11/01-marginal-reasoning/ (WDI life expectancy and
% health spending, 2019). Ledger: Arkes and Blumer (1985), Experiment 2. Brenland tax bands invented
% (code/shared/country TAX_BANDS); the bakery of the problem is invented in the problem itself.
\chapter{Marginal Reasoning}\label{ch:g11:marginal-reasoning}

In 2019, a country that spent 100 dollars a year on the health of each person could expect, from the line that best fits 193 countries, a life expectancy
about 3.4 years longer if it spent 100 dollars more. A country that spent 1000 dollars could expect about 0.46 years more from the same extra 100 dollars,
and one that spent 5000, about a month. The first dollars buy vaccines, clean births and antibiotics; the later ones buy rarer treatments for rarer illnesses.
% figdata: life_fit (n 193; gain100_at100 3.38, gain100_at1000 0.46, gain100_at5000 0.10 years)

The line says nothing about what causes what: richer countries both spend more on health and have cleaner water, better food and safer work. What it shows
is the shape that runs through this chapter. What matters for a decision is not the total that a choice brings, but what the next step brings and what it
costs. Economists call the next step the margin.

\begin{openingchart}{Life expectancy at birth against health spending per person, 193 countries, 2019. Spending is in dollars a year at purchasing-power
parity, on a times-ten scale. The curve is the line, on this scale, that passes closest to the points; the years gained from 100 more dollars fall as spending
rises.}{World Bank, World Development Indicators: life expectancy at birth, total; current health expenditure per capita, PPP}
\begin{tikzpicture}
\begin{semilogxaxis}[omscatter, log basis x=10, log ticks with fixed point, xtick={10,100,1000,10000}, width=0.84\linewidth, height=6.6cm,
  xmin=20, xmax=12000, ymin=50, ymax=90, xlabel={health spending per person (dollars a year)}, ylabel={life expectancy (years)}]
\addplot[omDef, only marks, mark=*, mark size=1.2pt, fill opacity=0.7]
  table[x=health, y=life] {figdata/grade-11/01-marginal-reasoning/out/life.csv};
\addplot[black, thick, no markers, sharp plot] table[x=health, y=life] {figdata/grade-11/01-marginal-reasoning/out/life_curve.csv};
\end{semilogxaxis}
\end{tikzpicture}
\end{openingchart}

\begin{recall}
From the grades 6--9 volume: the \emph{opportunity cost} of a choice is the best option given up; a \emph{sunk cost} is a cost already paid that cannot be
recovered, and it should not enter a decision; a firm's \emph{profit} is its revenue minus its costs; the \emph{marginal tax rate} is the share of one more
dalen of income taken in tax, and the \emph{average tax rate} is total tax divided by income. From mathematics: the derivative $f'(x)$ as the slope of $f$ at
$x$; a differentiable function on an interval has its maximum where $f' = 0$ if $f'$ goes from positive to negative there.
\end{recall}

\section{The next unit}\label{sec:g11:marginal-reasoning:next-unit}

\begin{definition}[Marginal quantity]\label{def:g11:marginal-reasoning:marginal-quantity}
If a total $F(q)$ depends on a quantity $q$, the \emph{marginal quantity}\index{marginal quantity} of $F$ at $q$ is the change in $F$ brought by one more unit
of $q$. When $F$ is differentiable and one unit is small, it is the derivative $F'(q)$.
\end{definition}

\begin{definition}[Marginal benefit, cost and revenue]\label{def:g11:marginal-reasoning:marginal-benefit}
The \emph{marginal benefit}\index{marginal benefit} of an activity is the extra benefit from one more unit of it; its \emph{marginal cost}\index{marginal cost}%
\phantomsection\label{def:g11:marginal-reasoning:marginal-cost} is the extra cost of that unit; a firm's \emph{marginal revenue}\index{marginal revenue}%
\phantomsection\label{def:g11:marginal-reasoning:marginal-revenue} is the extra revenue from selling one more unit.
\end{definition}

\begin{example}[A step and a slope]\label{ex:g11:marginal-reasoning:step}
A workshop's cost of making $q$ chairs a week is $C(q) = 200 + 10q + 0.5q^2$ dalens. Going from 20 to 21 chairs costs $C(21) - C(20) = 10 + 0.5 \times (441 - 400)
= 30.5$ dalens. The derivative $C'(q) = 10 + q$ gives $C'(20) = 30$. The two differ by half a dalen: the derivative is the slope at 20, the step is the average
slope between 20 and 21. When one unit is small compared with the quantity, the two are close, and the derivative is easier to work with.
\end{example}

The fixed 200 dalens a week (the rent of the workshop, say) appear in $C$ but not in $C'$: they do not change with the number of chairs, so they are no part of the
cost of any one more chair.

\section{Marginal benefit equals marginal cost}\label{sec:g11:marginal-reasoning:rule}

\begin{definition}[Net benefit]\label{def:g11:marginal-reasoning:net-benefit}
The \emph{net benefit}\index{net benefit} of doing $q$ units of an activity is its total benefit minus its total cost, $N(q) = B(q) - C(q)$.
\end{definition}

\begin{proposition}[The marginal rule]\label{prop:g11:marginal-reasoning:rule}
Let $B$ and $C$ be differentiable on an interval, and suppose that the \omterm{def:g11:marginal-reasoning:marginal-benefit}{marginal benefit} $B'$ minus the \omterm{def:g11:marginal-reasoning:marginal-benefit}{marginal cost} $C'$ is positive below $q^*$, zero at
$q^*$ and negative above it. Then the \omterm{def:g11:marginal-reasoning:net-benefit}{net benefit} is largest at $q^*$, where \omterm{def:g11:marginal-reasoning:marginal-benefit}{marginal benefit} equals \omterm{def:g11:marginal-reasoning:marginal-benefit}{marginal cost}.
\end{proposition}
\begin{proof}
$N'(q) = B'(q) - C'(q)$. It is positive below $q^*$, so $N$ increases up to $q^*$; it is negative above $q^*$, so $N$ decreases after it. Hence $N(q^*) \geq N(q)$
for every $q$ in the interval.
\end{proof}

In words: keep going while the next unit brings more than it costs, and stop where it brings just what it costs. The proposition needs the \omterm{def:g11:marginal-reasoning:marginal-benefit}{marginal benefit}
to fall, or the \omterm{def:g11:marginal-reasoning:marginal-benefit}{marginal cost} to rise, as $q$ grows: if the next unit always brought more than it cost, there would be no best stopping point inside the interval.

\begin{omfigure}
\begin{tikzpicture}[econ, x=0.42cm, y=0.055cm]
\begin{scope}
\draw[econ axes] (0,78) node[left, font=\small] {dalens} |- (11,0) node[below left=12pt and 0pt, font=\small] {$q$};
\draw[demand, domain=0:9, samples=37] plot (\x, {12*\x - 0.5*\x*\x}) node[right, font=\scriptsize] {$B$};
\draw[supply, domain=0:9, samples=37] plot (\x, {2*\x + 0.5*\x*\x}) node[right, font=\scriptsize] {$C$};
\draw[guide] (5,22.5) -- (5,0);
\draw[omThm, thick, <->] (5,23.5) -- (5,46.5);
\node[font=\scriptsize, left] at (4.9,35) {$N = 25$};
\foreach \y in {20,40,60} \draw (0,\y) -- (-0.2,\y) node[left, font=\footnotesize] {\y};
\node[below, font=\footnotesize] at (5,0) {5};
\node[below, font=\footnotesize] at (10,0) {10};
\end{scope}
\begin{scope}[xshift=7.2cm, y=0.3cm]
\draw[econ axes] (0,14) node[left, font=\small] {dalens} |- (11,0) node[below left=12pt and 0pt, font=\small] {$q$};
\draw[demand] (0,12) -- (10,2) node[right, font=\scriptsize] {$B'$};
\draw[supply] (0,2) -- (10,12) node[right, font=\scriptsize] {$C'$};
\draw[guide] (5,7) -- (5,0);
\node[eqpoint] at (5,7) {};
\foreach \y in {2,7,12} \draw (0,\y) -- (-0.2,\y) node[left, font=\footnotesize] {\y};
\node[below, font=\footnotesize] at (5,0) {5};
\node[below, font=\footnotesize] at (10,0) {10};
\end{scope}
\end{tikzpicture}

{\small Left: total benefit $B(q) = 12q - 0.5q^2$ and total cost $C(q) = 2q + 0.5q^2$. The gap between them, the \omterm{def:g11:marginal-reasoning:net-benefit}{net benefit}, is largest at $q = 5$, where the two
curves are equally steep. Right: the same choice seen at the margin: $B'(q) = 12 - q$ falls, $C'(q) = 2 + q$ rises, and they cross at $q = 5$.}
\end{omfigure}

\begin{model}{The optimising decision-maker}
\assumes A person or a firm chooses a quantity to make its \omterm{def:g11:marginal-reasoning:net-benefit}{net benefit} as large as possible, knowing its benefits and costs.
\predicts The chosen quantity sits where \omterm{def:g11:marginal-reasoning:marginal-benefit}{marginal benefit} equals \omterm{def:g11:marginal-reasoning:marginal-benefit}{marginal cost}; it moves when a \omterm{def:g11:marginal-reasoning:marginal-benefit}{marginal benefit} or a \omterm{def:g11:marginal-reasoning:marginal-benefit}{marginal cost} changes; it does not move when
  a cost that does not depend on the quantity changes, nor when a sunk cost changes.
\tested Consistent with it: producers expand when the price of their output rises and its cost does not, as American crude oil output rose by
  three-quarters within six years of high prices (\cref{ch:g10:supply}), and buyers of fuel cut their use when its price rises
  (\cref{ch:g10:elasticities}). A rule of thumb predicts the same directions; a test of the marginal rule proper is its claim that sunk costs do
  not matter, and the randomised test below finds that they do for a while.
\breaks Sunk costs do enter some decisions. In a theatre's 1982--83 season, the first buyers of a season ticket were given, at random, the normal price of 15
  dollars, 2 dollars off or 7 dollars off. Over the first five plays, buyers at the normal price used 4.11 tickets each, against 3.32 and 3.29 for the two
  discounted groups; over the last five plays the groups no longer differed (Arkes and Blumer 1985). The rule also needs the quantity to be divisible: when it
  comes in whole units, the slope can point to a quantity that cannot be bought.
\end{model}
% ledger: g11:arkes-blumer-1985

\section{Average and marginal}\label{sec:g11:marginal-reasoning:average}

The average cost of $q$ units is $A(q) = C(q)/q$. Average and marginal answer different questions: what each unit has cost so far, and what the next unit will
cost.

\begin{proposition}[The marginal pulls the average]\label{prop:g11:marginal-reasoning:average}
Let $C$ be differentiable and $q > 0$, with $A(q) = C(q)/q$ and $M(q) = C'(q)$. Then $A'(q) = \big(M(q) - A(q)\big)/q$: the average rises where the marginal is
above it, falls where the marginal is below it, and is flat where the two are equal.
\end{proposition}
\begin{proof}
By the quotient rule, $A'(q) = \big(C'(q)\,q - C(q)\big)/q^2 = \big(C'(q) - C(q)/q\big)/q = \big(M(q) - A(q)\big)/q$. Since $q > 0$, $A'(q)$ has the sign of $M(q) - A(q)$.
\end{proof}

The same holds for any average: a student whose next mark is above her average raises it. With $C(q) = 64 + q^2$, the \omterm{def:g11:marginal-reasoning:marginal-benefit}{marginal cost} is $2q$ and the average cost
$64/q + q$; they are equal at $q = 8$, where both are 16, and that is the bottom of the average cost curve.

\begin{omfigure}
\begin{tikzpicture}[econ, x=0.42cm, y=0.12cm]
\draw[econ axes] (0,36) node[left, font=\small] {dalens} |- (17,0) node[below left=12pt and 0pt, font=\small] {$q$};
\draw[supply, domain=2:16, samples=50] plot (\x, {64/\x + \x}) node[right, font=\scriptsize] {average cost $64/q + q$};
\draw[demand, domain=0:16] plot (\x, {2*\x}) node[right, font=\scriptsize] {marginal cost $2q$};
\draw[guide] (8,16) -- (8,0);
\node[eqpoint] at (8,16) {};
\foreach \y in {16,32} \draw (0,\y) -- (-0.2,\y) node[left, font=\footnotesize] {\y};
\node[below, font=\footnotesize] at (8,0) {8};
\node[below, font=\footnotesize] at (16,0) {16};
\end{tikzpicture}

{\small Average cost falls while \omterm{def:g11:marginal-reasoning:marginal-benefit}{marginal cost} lies below it and rises once \omterm{def:g11:marginal-reasoning:marginal-benefit}{marginal cost} is above it; the \omterm{def:g11:marginal-reasoning:marginal-benefit}{marginal cost} line crosses the average cost curve at its
lowest point, $q = 8$.}
\end{omfigure}

Taxes on income show the two side by side. In Brenland, the first \dal{12000} of a year's income are not taxed; income from \dal{12000} to \dal{40000} is taxed at
20\,\%, from \dal{40000} to \dal{100000} at 35\,\%, and above \dal{100000} at 45\,\%. A person earning \dal{50000} pays $0.20 \times 28000 + 0.35 \times 10000 = 9100$ dalens:
an average rate of 18.2\,\%, a marginal rate of 35\,\%.
% figdata: tax_rates (50000: marginal 35, average 18.2); code/shared/country TAX_BANDS

\begin{omchart}{Brenland's income tax: the marginal rate, in steps, and the average rate, against income in dalens a year.}{Brenland tax law (invented)}
\begin{tikzpicture}
\begin{axis}[omtime, width=0.84\linewidth, height=5.4cm, xmin=0, xmax=150000, ymin=0, ymax=50, ylabel={percent of income},
  xlabel={income (dalens a year)}, scaled x ticks=false, xticklabel style={/pgf/number format/fixed, /pgf/number format/1000 sep={\,}},
  xtick={0,40000,100000,150000}, legend pos=south east, legend style={font=\scriptsize}]
\addplot[omThm, thick, const plot mark left] table[x=income, y=marginal] {figdata/grade-11/01-marginal-reasoning/out/tax_rates.csv};
\addlegendentry{marginal rate}
\addplot[omDef, thick] table[x=income, y=average] {figdata/grade-11/01-marginal-reasoning/out/tax_rates.csv}; \addlegendentry{average rate}
\end{axis}
\end{tikzpicture}
\end{omchart}

\begin{measure}[Marginal and average tax rates]
The average rate says how much of a person's income the tax takes; the marginal rate says how much of one more dalen it takes, and it is the one that matters for
a decision to work an extra hour or take a better-paid job. The average rate is never above the marginal rate in a schedule like Brenland's, and it climbs towards
the top rate without reaching it. Measuring the marginal rate a person actually faces means adding every tax and every benefit that changes with income: social
contributions, and benefits that are cut as earnings rise, can push it well above the rate printed in the tax law.
\end{measure}

\section{What is gone is gone}\label{sec:g11:marginal-reasoning:sunk}

A sunk cost is the same whatever is chosen now, so it cancels out of every comparison: it shifts the \omterm{def:g11:marginal-reasoning:net-benefit}{net benefit} of every option by the same amount and leaves the
best option unchanged. A firm that spent \dal{80000} developing a product and needs \dal{30000} more to finish it should finish if the product will bring in more than
\dal{30000}, and stop otherwise. The \dal{80000} decide nothing. What they can do is teach: a firm that keeps losing such sums learns something about how it picks its
projects.

\section{Exercises}

\begin{exercise}[$\star$]\label{exo:g11:marginal-reasoning:1}
The total cost of making 0 to 5 cakes is 20, 26, 30, 36, 46 and 60 dalens. Compute the \omterm{def:g11:marginal-reasoning:marginal-benefit}{marginal cost} of each cake.
\end{exercise}

\begin{exercise}[$\star$]\label{exo:g11:marginal-reasoning:2}
A student's total benefit from 0 to 5 hours of revision, in points on the exam, is 0, 30, 50, 62, 68 and 70. Each hour of revision costs her the equivalent of 10
points of other work. Compute the \omterm{def:g11:marginal-reasoning:marginal-benefit}{marginal benefit} of each hour, and the number of hours she should revise.
\end{exercise}

\begin{exercise}[$\star$]\label{exo:g11:marginal-reasoning:3}
A cost function is $C(q) = 50 + 4q + 0.5q^2$. Compute the \omterm{def:g11:marginal-reasoning:marginal-benefit}{marginal cost} $C'(10)$, and compare it with $C(11) - C(10)$.
\end{exercise}

\begin{exercise}[$\star$]\label{exo:g11:marginal-reasoning:4}
Read the opening chart. About what life expectancy goes with spending of 100 dollars a person? Of 1000? Of 10000? By about how many years does it rise each time
spending is multiplied by ten?
\end{exercise}

\begin{exercise}[$\star$]\label{exo:g11:marginal-reasoning:5}
A town has spent 2 million dalens on half a bridge. Finishing it would cost 1.5 million dalens more, and the finished bridge would bring benefits worth 1.2
million dalens. Should the town finish it? Would your answer change if it had spent 3 million dalens so far?
\end{exercise}

\begin{exercise}[$\star\star$]\label{exo:g11:marginal-reasoning:6}
The \omterm{def:g11:marginal-reasoning:net-benefit}{net benefit} of $q$ units is $N(q) = 40q - 2q^2 - 50$. Find the quantity that makes it largest, and the largest \omterm{def:g11:marginal-reasoning:net-benefit}{net benefit}.
\end{exercise}

\begin{exercise}[$\star\star$]\label{exo:g11:marginal-reasoning:7}
The benefit of $q$ units is $B(q) = 30q - q^2$ and the cost $C(q) = 6q + 0.5q^2$. Find the \omterm{def:g11:marginal-reasoning:marginal-benefit}{marginal benefit} and \omterm{def:g11:marginal-reasoning:marginal-benefit}{marginal cost}, the best quantity, and the \omterm{def:g11:marginal-reasoning:net-benefit}{net
benefit} there.
\end{exercise}

\begin{exercise}[$\star\star$]\label{exo:g11:marginal-reasoning:8}
With Brenland's tax bands, compute the tax, the average rate and the marginal rate for an income of \dal{30000}.
\end{exercise}

\begin{exercise}[$\star\star$]\label{exo:g11:marginal-reasoning:9}
Same question for an income of \dal{120000}. How much of an extra \dal{1000} does this person keep?
\end{exercise}

\begin{exercise}[$\star\star$]\label{exo:g11:marginal-reasoning:10}
In \cref{exo:g11:marginal-reasoning:7}, a charge of 3 dalens is added to each unit. What happens to the best quantity? What if, instead, a single charge of 50
dalens is added whatever the quantity?
\end{exercise}

\begin{exercise}[$\star\star$]\label{exo:g11:marginal-reasoning:11}
A cost function is $C(q) = 100 + q^2$. Compute the average and the \omterm{def:g11:marginal-reasoning:marginal-benefit}{marginal cost} at $q = 5$, $10$ and $15$, and say whether the average is rising or falling at each.
\end{exercise}

\begin{exercise}[$\star\star\star$]\label{exo:g11:marginal-reasoning:12}
In the theatre experiment of the model box, what would the optimising decision-maker predict for the three groups? Which part of the result breaks the prediction,
and which part fits it? Why does it matter that the discounts were handed out at random?
\end{exercise}

\begin{exercise}[$\star\star\star$]\label{exo:g11:marginal-reasoning:13}
A list of $n$ numbers has mean $m$. Show that adding a number $x$ raises the mean if $x > m$ and lowers it if $x < m$, and compute the new mean. Relate this to
\cref{prop:g11:marginal-reasoning:average}.
\end{exercise}

\begin{exercise}[$\star\star\star$]\label{exo:g11:marginal-reasoning:14}
A firm can run only a whole number $q$ of delivery vans, and its \omterm{def:g11:marginal-reasoning:net-benefit}{net benefit} is $N(q) = 20.3q - q^3/3$ dalens a day. Where is $N'(q) = 0$? Compute $N(4)$ and
$N(5)$. Which number of vans should the firm run, and why does rounding the answer of the derivative mislead here?
\end{exercise}

\begin{exercise}[$\star\star\star$]\label{exo:g11:marginal-reasoning:15}
A Brenland benefit pays \dal{6000} a year to a person with no earnings, and is cut by 50 centimes for every dalen earned until it reaches zero. Using Brenland's
tax bands, compute the person's total income (earnings plus benefit minus tax) at earnings of \dal{0}, \dal{10000} and \dal{14000}. What is the marginal rate on the
next dalen earned at \dal{10000}? Compare it with the top tax rate.
\end{exercise}

\section{Problem: A Bakery at the Margin}

\begin{problem}[{Weekend problem --- a bakery chooses how many loaves to bake}]
\label{pb:g11:marginal-reasoning:1}
A bakery in a town in Brenland sells its own loaves. If it bakes and sells $q$ loaves a day, the price at which they all sell is $p = 6 - 0.01q$ dalens. Its cost
is $C(q) = 120 + 0.9q + 0.005q^2$ dalens a day: 120 dalens of daily rent, and flour, power and wages that grow with the number of loaves. Last year the baker
bought an oven for \dal{9000}; it cannot be resold.

\medskip
\noindent\textbf{Part I --- Revenue and cost at the margin.}
\begin{enumerate}
\item Write the bakery's revenue $R(q)$, and compute the price and the revenue at $q = 100$.
\item Compute the \omterm{def:g11:marginal-reasoning:marginal-benefit}{marginal revenue} $R'(q)$, its value at $q = 100$, and compare with $R(101) - R(100)$.
\item Compute the \omterm{def:g11:marginal-reasoning:marginal-benefit}{marginal cost} $C'(q)$ and its value at $q = 100$.
\item Compute the average cost at $q = 100$. Why does it differ from the \omterm{def:g11:marginal-reasoning:marginal-benefit}{marginal cost}?
\item Should the bakery bake a 101st loaf? Answer with marginal quantities only.
\end{enumerate}

\medskip
\noindent\textbf{Part II --- The best number of loaves.}
\begin{enumerate}[resume]
\item Find the number of loaves at which \omterm{def:g11:marginal-reasoning:marginal-benefit}{marginal revenue} equals \omterm{def:g11:marginal-reasoning:marginal-benefit}{marginal cost}.
\item Compute the price at that quantity.
\item Compute the profit.
\item Check that the profit is lower at 169 and at 171 loaves.
\item Compute the average and \omterm{def:g11:marginal-reasoning:marginal-benefit}{marginal cost} at the best quantity. Is the average cost rising or falling there?
\end{enumerate}

\medskip
\noindent\textbf{Part III --- Flour costs more.}
\begin{enumerate}[resume]
\item The price of flour rises, adding 0.60 dalens to the cost of each loaf. Write the new cost and \omterm{def:g11:marginal-reasoning:marginal-benefit}{marginal cost}.
\item Find the new best quantity.
\item Compute the new price.
\item Compute the new profit.
\item The cost of each loaf rose by 0.60 dalens. By how much did the price rise? Who bears the rest?
\end{enumerate}

\medskip
\noindent\textbf{Part IV --- Costs that do not move the margin.}
\begin{enumerate}[resume]
\item Does the oven's \dal{9000} change the best number of loaves? Explain.
\item The baker says: ``We must sell more loaves to pay off the oven.'' What is wrong with this?
\item After the flour rise, the landlord raises the rent to 250 dalens a day. What happens to the best quantity and to the profit?
\item The lease can be ended at any time. Up to what daily rent should the bakery stay open?
\item State the result of the problem: the profit-maximising number of loaves and the profit, before and after the flour price rise.
\end{enumerate}
\end{problem}
